% GENERATED FILE. Edit canonical lessons, then run npm run generate:books.
% canonical-source-hash: fnv1a64:28306bebbbb80a65
% canonical-lessons: SA-C542-svikarah, SA-C542-avadharayati, SA-C542-utkrstah, SA-C542-rucirah, SA-C542-tatasthah

\chapter{Making Up One's Mind}
\label{ch:sa-making-up-one-s-mind}
\hlchaptermodality{542}

\begin{chapteropening}
\textbf{By the end of this chapter:} \emph{I can say what I will choose and why.}

Complete the last lesson of chapter 542: I can say what I will choose and why.
\end{chapteropening}

\section[\texorpdfstring{\sk{स्वीकारः} (svīkāraḥ)}{स्वीकारः (svīkāraḥ)}]{\sk{स्वीकारः} (svīkāraḥ) — acceptance}

\label{lesson:SA-C542-svikarah}



\begin{quote}
Before the new one: say the Sanskrit for the other one, then the Sanskrit for different.
\end{quote}

\subsection*{You'll want to know: \sk{स्वीकारः}}
\textbf{\sk{स्वीकारः}} — \emph{svīkāraḥ} — \textquotedblleft{}acceptance\textquotedblright{}.

\begin{quote}
\sk{यदि} \sk{भवतः} \sk{स्वीकारः} \sk{अस्ति}, \sk{तर्हि} \sk{वयं} \sk{कार्यम्} \sk{आरभामहे।}
\end{quote}

\begin{quote}
\emph{yadi bhavataḥ svīkāraḥ asti, tarhi vayaṃ kāryam ārabhāmahe.}
\end{quote}

\begin{quote}
If you accept, we will begin the work.
\end{quote}

A masculine noun.

\begin{cousinweb}
\textbf{\sk{स्वीकारः}} (\emph{svīkāraḥ}) is built from \emph{sva}, \textquotedblleft{}one's own\textquotedblright{}, in the form \emph{svī}, and \emph{kāra}, \textquotedblleft{}making\textquotedblright{}, from the root √kṛ; MW gives \emph{svīkāra} as making one's own, and reception (MW).
\end{cousinweb}

\begin{grammarlens}[title={What you've built}]
One more word for saying what happens if.
\end{grammarlens}

\subsection*{Your turn}
\begin{itemize}
\raggedright
  \item \emph{Say it:} \emph{svīkāraḥ}
  \item \emph{Say it:} \emph{svīkāraḥ}, once more
  \item \emph{Say it:} \emph{yadi bhavataḥ svīkāraḥ asti, tarhi vayaṃ kāryam ārabhāmahe.} — If you accept, we will begin the work.
  \item \emph{Recall:} say the Sanskrit for the other one, then the Sanskrit for different, then say \emph{svīkāraḥ} again
\end{itemize}

\begin{tcolorbox}[breakable,colback=teal!4,colframe=teal!35!black,title={Before you move on}]
What does \sk{स्वीकारः} mean? (\textquotedblleft{}Acceptance\textquotedblright{}.) Say it once more.
\end{tcolorbox}
\section[\texorpdfstring{\sk{अवधारयति} (avadhārayati)}{अवधारयति (avadhārayati)}]{\sk{अवधारयति} (avadhārayati) — makes sure}

\label{lesson:SA-C542-avadharayati}



\begin{quote}
Before the new one: say the Sanskrit for different, then the Sanskrit for acceptance.
\end{quote}

\subsection*{You'll want to know: \sk{अवधारयति}}
\textbf{\sk{अवधारयति}} — \emph{avadhārayati} — \textquotedblleft{}makes sure\textquotedblright{}.

\begin{quote}
\sk{यदि} \sk{भवान्} \sk{मार्गं} \sk{न} \sk{जानाति}, \sk{तर्हि} \sk{पूर्वम्} \sk{अवधारयतु।}
\end{quote}

\begin{quote}
\emph{yadi bhavān mārgaṃ na jānāti, tarhi pūrvam avadhārayatu.}
\end{quote}

\begin{quote}
If you don't know the way, make sure of it first.
\end{quote}

Past (he …): \sk{सः} \sk{मार्गम्} \sk{अवधारितवान्}, \emph{saḥ mārgam avadhāritavān}.

\begin{cousinweb}
\textbf{\sk{अवधारयति}} (\emph{avadhārayati}) is built on \emph{ava}, \textquotedblleft{}down\textquotedblright{}, and the causative of the root √dhṛ, \textquotedblleft{}to hold\textquotedblright{}, which also gives \emph{dhārayati}; MW gives \emph{avadhārayati} as to consider, to ascertain, to determine accurately (MW).
\end{cousinweb}

\begin{grammarlens}[title={What you've built}]
One more word for saying what happens if.
\end{grammarlens}

\subsection*{Your turn}
\begin{itemize}
\raggedright
  \item \emph{Say it:} \emph{avadhārayati}
  \item \emph{Say it:} \emph{avadhārayati}, once more
  \item \emph{Say it:} \emph{yadi bhavān mārgaṃ na jānāti, tarhi pūrvam avadhārayatu.} — If you don't know the way, make sure of it first.
  \item \emph{Recall:} say the Sanskrit for different, then the Sanskrit for acceptance, then say \emph{avadhārayati} again
\end{itemize}

\begin{tcolorbox}[breakable,colback=teal!4,colframe=teal!35!black,title={Before you move on}]
What does \sk{अवधारयति} mean? (\textquotedblleft{}Makes sure\textquotedblright{}.) Say it once more.
\end{tcolorbox}
\section[\texorpdfstring{\sk{उत्कृष्टः} (utkṛṣṭaḥ)}{उत्कृष्टः (utkṛṣṭaḥ)}]{\sk{उत्कृष्टः} (utkṛṣṭaḥ) — excellent}

\label{lesson:SA-C542-utkrstah}



\begin{quote}
Before the new one: say the Sanskrit for acceptance, then the Sanskrit for makes sure.
\end{quote}

\subsection*{You'll want to know: \sk{उत्कृष्टः}}
\textbf{\sk{उत्कृष्टः}} — \emph{utkṛṣṭaḥ} — \textquotedblleft{}excellent\textquotedblright{}.

\begin{quote}
\sk{यदि} \sk{एतत्} \sk{वस्त्रम्} \sk{उत्कृष्टम्} \sk{अस्ति}, \sk{तर्हि} \sk{अहं} \sk{तत्} \sk{क्रीणामि।}
\end{quote}

\begin{quote}
\emph{yadi etat vastram utkṛṣṭam asti, tarhi ahaṃ tat krīṇāmi.}
\end{quote}

\begin{quote}
If this cloth is excellent, I'll buy it.
\end{quote}

\begin{cousinweb}
\textbf{\sk{उत्कृष्टः}} (\emph{utkṛṣṭaḥ}) is built on \emph{ud}, \textquotedblleft{}up\textquotedblright{}, and the root √kṛṣ, \textquotedblleft{}to draw\textquotedblright{}, which also gives \emph{karṣati}, as its past participle: drawn up high; MW gives \emph{utkṛṣṭa} as excellent, eminent (MW).
\end{cousinweb}

\begin{grammarlens}[title={What you've built}]
One more word for saying what happens if.
\end{grammarlens}

\subsection*{Your turn}
\begin{itemize}
\raggedright
  \item \emph{Say it:} \emph{utkṛṣṭaḥ}
  \item \emph{Say it:} \emph{utkṛṣṭaḥ}, once more
  \item \emph{Say it:} \emph{yadi etat vastram utkṛṣṭam asti, tarhi ahaṃ tat krīṇāmi.} — If this cloth is excellent, I'll buy it.
  \item \emph{Recall:} say the Sanskrit for acceptance, then the Sanskrit for makes sure, then say \emph{utkṛṣṭaḥ} again
\end{itemize}

\begin{tcolorbox}[breakable,colback=teal!4,colframe=teal!35!black,title={Before you move on}]
What does \sk{उत्कृष्टः} mean? (\textquotedblleft{}Excellent\textquotedblright{}.) Say it once more.
\end{tcolorbox}
\section[\texorpdfstring{\sk{रुचिरः} (ruciraḥ)}{रुचिरः (ruciraḥ)}]{\sk{रुचिरः} (ruciraḥ) — pleasing; lovely}

\label{lesson:SA-C542-rucirah}



\begin{quote}
Before the new one: say the Sanskrit for makes sure, then the Sanskrit for excellent.
\end{quote}

\subsection*{You'll want to know: \sk{रुचिरः}}
\textbf{\sk{रुचिरः}} — \emph{ruciraḥ} — \textquotedblleft{}pleasing; lovely\textquotedblright{}.

\begin{quote}
\sk{यदि} \sk{एतत्} \sk{पुष्पं} \sk{रुचिरम्} \sk{अस्ति}, \sk{तर्हि} \sk{अहं} \sk{तत्} \sk{चिनोमि।}
\end{quote}

\begin{quote}
\emph{yadi etat puṣpaṃ ruciram asti, tarhi ahaṃ tat cinomi.}
\end{quote}

\begin{quote}
If this flower is lovely, I'll pick it.
\end{quote}

\begin{cousinweb}
\textbf{\sk{रुचिरः}} (\emph{ruciraḥ}) comes from the root √ruc, \textquotedblleft{}to shine, to please\textquotedblright{}, which also gives \emph{rocate}; MW gives \emph{rucira} as bright, beautiful, and pleasant, liked by someone (MW).
\end{cousinweb}

\begin{grammarlens}[title={What you've built}]
One more word for saying what happens if.
\end{grammarlens}

\subsection*{Your turn}
\begin{itemize}
\raggedright
  \item \emph{Say it:} \emph{ruciraḥ}
  \item \emph{Say it:} \emph{ruciraḥ}, once more
  \item \emph{Say it:} \emph{yadi etat puṣpaṃ ruciram asti, tarhi ahaṃ tat cinomi.} — If this flower is lovely, I'll pick it.
  \item \emph{Recall:} say the Sanskrit for makes sure, then the Sanskrit for excellent, then say \emph{ruciraḥ} again
\end{itemize}

\begin{tcolorbox}[breakable,colback=teal!4,colframe=teal!35!black,title={Before you move on}]
What does \sk{रुचिरः} mean? (\textquotedblleft{}Pleasing; lovely\textquotedblright{}.) Say it once more.
\end{tcolorbox}
\section[\texorpdfstring{\sk{तटस्थः} (taṭasthaḥ)}{तटस्थः (taṭasthaḥ)}]{\sk{तटस्थः} (taṭasthaḥ) — neutral}

\label{lesson:SA-C542-tatasthah}



\begin{quote}
Before the new one: say the Sanskrit for excellent, then the Sanskrit for pleasing.
\end{quote}

\subsection*{You'll want to know: \sk{तटस्थः}}
\textbf{\sk{तटस्थः}} — \emph{taṭasthaḥ} — \textquotedblleft{}neutral\textquotedblright{}.

\begin{quote}
\sk{यदि} \sk{भवान्} \sk{तटस्थः} \sk{अस्ति}, \sk{तर्हि} \sk{भवान्} \sk{निर्णयं} \sk{करोतु।}
\end{quote}

\begin{quote}
\emph{yadi bhavān taṭasthaḥ asti, tarhi bhavān nirṇayaṃ karotu.}
\end{quote}

\begin{quote}
If you are neutral, you decide.
\end{quote}

\begin{cousinweb}
\textbf{\sk{तटस्थः}} (\emph{taṭasthaḥ}) is built from \emph{taṭa}, \textquotedblleft{}a bank, a slope\textquotedblright{}, and \emph{stha}, \textquotedblleft{}standing\textquotedblright{}: standing on the bank, out of the stream; MW gives \emph{taṭastha} also as an indifferent person, neither friend nor foe (MW).
\end{cousinweb}

\begin{grammarlens}[title={What you've built}]
One more word for saying what happens if.
\end{grammarlens}

\subsection*{Your turn}
\begin{itemize}
\raggedright
  \item \emph{Say it:} \emph{taṭasthaḥ}
  \item \emph{Say it:} \emph{taṭasthaḥ}, once more
  \item \emph{Say it:} \emph{yadi bhavān taṭasthaḥ asti, tarhi bhavān nirṇayaṃ karotu.} — If you are neutral, you decide.
  \item \emph{Recall:} say the Sanskrit for excellent, then the Sanskrit for pleasing, then say \emph{taṭasthaḥ} again
\end{itemize}

\begin{tcolorbox}[breakable,colback=teal!4,colframe=teal!35!black,title={Before you move on}]
What does \sk{तटस्थः} mean? (\textquotedblleft{}Neutral\textquotedblright{}.) Say it once more.
\end{tcolorbox}
